\subsection{UNIFORMIZABLE SPACES}

In Chapter II, § 4, no. 1, we posed the problem of characterizing uniformizable topological spaces. The solution is given by the following theorem:

THEOREM 2. A topological space X is uniformizable if and only if it satisfies the following axiom:

\((\mathrm{O}_{\mathrm{IV}})\) Given any point \(x_0 \in X\) and any neighbourhood \(V\) of \(x_0\) , there exists a continuous real-valued function on \(X\) which takes its values in \([0,1]\) , is equal to 0 at \(x_0\) , and is equal to 1 on CV.

The condition is necessary. For if there is a uniformity on X compatible with the topology of X, then by Theorem 1 this uniformity can be defined by a family \((f_{t})\) of pseudometrics on X, and we may assume with no loss of generality that this family is saturated (no. 2). From the definition of the entourages of the uniformity defined by such a family of pseudometrics, there is a pseudometric \(f_{\alpha}\) of the family \((f_{t})\) , and a number \(a > 0\) , such that \(f_{\alpha}(x_0, x) \geqslant a\) for all \(x \in \mathbb{C}\mathrm{V}\) . It follows that the function \(g(x) = \inf \left(1, \frac{1}{a} f_{\alpha}(x_0, x)\right)\) satisfies all the conditions laid down in \((\mathrm{O}_{\mathrm{IV}})\) .

The condition is sufficient. For let \(\Phi\) be the set of all continuous mappings of X into {[}o, i{]}. Axiom (O \(_{IV}\) ) shows that the coarsest uniformity with respect to which all the functions belonging to \(\Phi\) are uniformly continuous is compatible with the topology of X (Chapter II, § 2, no. 3).

DEFINITION 4. A topological space is said to be completely regular if it is uniformizable and Hausdorff.

Equivalently, in view of Theorem 2, a space is completely regular if it satisfies axioms (H) {[}cf.~Chapter I, § 8, no. 1, Proposition 1{]} and \((\mathbf{O}_{\mathbf{IV}})\) .

Remark. Axiom \((\mathrm{O}_{\mathrm{IV}})\) implies \((\mathrm{O}_{\mathrm{III}})\) (cf.~Chapter I, § 8, no. 4), for if V is a neighbourhood of \(x_0\) and if \(f\) is a continuous real-valued function on X with values in {[}o, r{]}, such that \(f(x_0) = o\) , \(f(x) = 1\) for all \(x \in \complement V\) , then the set \(\overline{f}([0, 1/2])\) is a closed neighbourhood of \(x_0\) contained in V. In particular, every completely regular space is regular (which justifies the terminology). But there are examples of regular spaces which are not completely regular (*), so that \((\mathrm{O}_{\mathrm{III}})\) does not imply \((\mathrm{O}_{\mathrm{IV}})\) .

Every compact space is completely regular (Chapter II, § 4, no. 1, Theorem 1) and therefore so is every subspace of a compact space. We can now complete this proposition by proving its converse:

PROPOSITION 3. A topological space X is completely regular if and only if it is homeomorphic to a subspace of a compact space.

Consider the coarsest uniformity on X with respect to which all continuous mappings of X into {[}o, 1{]} are uniformly continuous; we have already used this uniformity in the proof of Theorem 2, where we saw that it is compatible with the topology of X if X is uniformizable. Furthermore, this uniformity is a structure of a precompact space, by the compactness of the interval {[}o, 1{]} and Proposition 3 of Chapter II, § 4, no. 2. If X is Hausdorff, the completion of X with respect to this is therefore compact, and the proposition is proved.

We can therefore say that a completely regular space can be embedded in a compact space. It is often convenient to present this result in the following way:

In general, a cube is a topological space \(K^{I}\) , the product of a family of topological spaces each identical with a compact interval K of R, indexed by an arbitrary set I. If I is finite and has n elements, we recover the notion of an n-dimensional closed cube, which was defined in Chapter VI, § 1, no. 1. A cube is a compact space (Chapter I, § 9, no. 5, Theorem 3).

PROPOSITION 4. If a topological space X is completely regular, it is homeomorphic to a subspace of a cube.

Let \((f_{i})_{i\in \mathbf{I}}\) denote the family of all continuous mappings of \(\mathbf{X}\) into \(\mathbf{K} = [\mathrm{o},\mathrm{i}]\) , and let \(g\) denote the mapping \(x\to (f_i(x))\) of \(\mathbf{X}\) into

\(K^{I}\) . If x, y are any two distinct points of X, it follows from axioms (H) and \((\mathrm{O}_{\mathrm{IV}})\) that there is an index t such that \(f_{t}(x) \neq f_{t}(y)\) , and therefore g is a one-to-one mapping of X into \(K^{I}\) . Moreover, it is immediate that g is an isomorphism of the coarsest uniformity on X for which all the \(f_{t}\) are uniformly continuous, onto the uniformity induced on \(g(\mathbf{X})\) by the product uniformity of \(K^{I}\) ; a fortiori, g is a homeomorphism of X onto \(g(\mathbf{X})\) .

